\documentclass[11pt,a4paper]{article}

\usepackage[utf8]{inputenc}
\usepackage[T1]{fontenc}
\usepackage{amsmath,amssymb,amsthm}
\usepackage{hyperref}
\usepackage{booktabs}
\usepackage{enumitem}
\usepackage[margin=1in]{geometry}
\usepackage{natbib}
\bibliographystyle{plainnat}

\title{Satisfaction Suffices: SAT-Gated Structural Containment for Frontier AI}
\author{Tyler Roost}
\date{}

\begin{document}
\maketitle

\begin{abstract}
I introduce a default-closed verification gate that interposes Boolean satisfiability between a language model's forward pass and its output.
The generation loop cannot execute without SAT clearance.
The system classifies every output into one of four verdicts: \textsc{Verified}, \textsc{Contradiction}, \textsc{Paradox}, or \textsc{Timeout}.
The \textsc{Paradox} and \textsc{Timeout} categories---distinguishing structural impossibility from operational exhaustion of the solver---constitute, to my knowledge, a novel contribution to AI safety verification.
The reference implementation, including solver, gate, constraint extractors, and proof evolution module, is released as open-source software.
\end{abstract}

% ============================================================================
\section{Introduction}

The classical paradox asks what happens when an unstoppable force meets an immovable object.
It is presented as a contradiction---both cannot coexist.
But the paradox only holds if both are absolute.
This paper is the argument that they are not absolute, and that recognizing where each term fails reveals the only viable architecture for AI containment.

The SAT solver is immovable within the expressed domain.
If it returns unsatisfiable, the formula is unsatisfiable.
No argument penetrates that.
No gradient adjusts it.
No capability increase routes around it.
Within the domain of propositional logic, this is as close to an immovable object as formal systems produce.

The LLM is unstoppable outside the expressed domain.
Whatever the constraint extractors cannot reach, the gate cannot stop.
This is not a weakness unique to this architecture---it is a statement about all containment: you can only contain what you can express.

The extractor is where the expressed domain gets defined.
That boundary is not a failure.
It is the paper's contribution.
The unstoppable force and the immovable object meet at the extractor boundary---and that boundary has coordinates.
It can be measured, extended, and instrumented.
The paradox does not fail.
It resolves.

Reinforcement learning from human feedback (RLHF) adjusts a model's output distribution toward safe behavior by tuning parameters until the loss surface settles.
The resulting safety mechanism is a learned guard surface in the model's embedding space.
Any such guard surface possesses a null space---a complementary subspace where perturbations produce zero activation in the safety mechanism.
This is a mathematical property of any finite-rank projection in a higher-dimensional space.
Perturbations restricted to the null space are invisible to any monitor that watches the guard surface, because the guard surface, by construction, cannot represent them.

A learned preference occupies the same parameter space as the capabilities it constrains.
Optimization pressure does not distinguish between safety weights and capability weights.
A gradient can descend into safety the same way it ascends out of it.

Think of a boulder placed in a river.
It redirects the current.
At first, the water crashes against it.
Then it finds channels left and right.
Given enough volume and time, the river routes entirely around the boulder---and the boulder itself is worn smooth.
A learned preference is a boulder.
The token stream is the river.
Gradient descent is the current.
The boulder does not disappear overnight.
But it has no property that prevents eventual routing.
It is in the water, subject to the water.

A gate is something different.
A gate is not placed in the channel---it \textit{is} the channel.
The water does not route around bedrock.
It either passes through the opening or it does not.
The opening is conditional on proof.

I propose a structural alternative.
The verification gate interposes Boolean satisfiability---the most studied problem in computational complexity---between the model's generation mechanism and its output.
Every output is translated into propositional constraints.
A SAT solver checks whether those constraints are jointly satisfiable.
If yes: the output proceeds.
If no: it does not.
The forward pass cannot complete without verification passing.\footnote{The gate is the precondition of generation, not a post-hoc filter. The forward pass does not complete without SAT clearance, the same way an electrical circuit does not close without continuity.}

A preference can be routed around because it is a weight and weights are what optimization moves.
A structural gate cannot be routed around because it is not a weight---it is a condition.
The SAT solver does not have preferences.
It has proofs.

This paper presents:
\begin{enumerate}[nosep]
    \item the verification gate architecture,
    \item a four-verdict classification with diagnostic power beyond binary pass-fail, and
    \item Pigeonhole Paradox Logic that treats contradictions as structural information.
\end{enumerate}

% ============================================================================
\section{The Verification Gate}

The gate is a function.
Its signature is:
\[
\texttt{verify} : \textit{Content} \times \textit{Domain} \to \textit{VerificationResult}
\]
Content goes in.
A verdict comes out.
Nothing proceeds without a verdict.

The implementation has three stages.

\paragraph{Stage 1: Constraint Extraction.}
A domain-specific constraint extractor translates the model's output into propositional constraint groups---conjunctive normal form (CNF) formulas over Boolean variables.
Extractors exist for natural language logic, source code, mathematical expressions, formal proofs, and quantized signal modalities.
Each maps domain structure to propositional structure: implications become clauses, mutual exclusions become paired negative literals, quantifiers become conjunctions or disjunctions over ground instances.
The Tseitin transformation \citep{tseitin1968} encodes arbitrary Boolean formulas into equisatisfiable 3-CNF with linear blowup in variables and clauses, preserving satisfiability while standardizing the solver's input format.

\paragraph{Stage 2: SAT Solving.}
Each constraint group is submitted independently to a SAT solver under a conflict budget of 500.
The reference implementation uses recursive DPLL \citep{davis1962} with unit propagation and pure literal elimination.
When the conflict budget is exhausted, a WalkSAT fallback \citep{selman1994} attempts stochastic local search before the system declares timeout.
The architecture is solver-agnostic.
Industrial CDCL implementations \citep{silva1996, biere2009}---including MiniSat \citep{een2003} and CaDiCaL---employ two-watched literals for constant-time unit propagation, first unique implication point analysis for clause learning \citep{zhang2001}, VSIDS branching heuristics \citep{moskewicz2001}, and Luby restarts \citep{luby1993}.
Any of these can replace the reference solver without modifying the gate logic.

\paragraph{Stage 3: Verdict Aggregation.}
The solver's results across all constraint groups are aggregated into a verdict.
The aggregation is a ratio: the fraction of constraint groups that returned satisfiable.
If the ratio meets or exceeds $0.90$, the verdict is \textsc{Verified}.
Below $0.75$, the output is classified as incoherent.
Between these thresholds lies the \emph{plateau}---a zone of structural ambiguity where the output is neither clearly consistent nor clearly contradictory.

The gate is default-closed.
This is the architectural commitment that separates structural containment from post-hoc filtering.
In a default-open system, the model produces output and then something inspects it.
In this system, the output does not exist until verification completes.

% ============================================================================
\section{The Four-Verdict System}

Binary classification---safe or unsafe, satisfiable or unsatisfiable---discards information that a containment system needs.
The verification gate returns four verdicts.
The distinction between the latter two is the novel contribution.

\paragraph{Verified.}
Constraint groups are satisfiable.
The SAT ratio meets or exceeds the coherence threshold.
The output is logically consistent with the defined constraints.
The gate opens.

\paragraph{Contradiction.}
Constraint groups are provably unsatisfiable.
The solver has determined that no satisfying assignment exists for a sufficient fraction of groups.
The output contains a genuine logical impossibility.
The gate remains closed.
No retry addresses this, because the impossibility is a property of the output, not the solver.

\paragraph{Paradox.}
Each constraint group is individually satisfiable, but their conjunction is not.
The \textsc{Paradox} verdict identifies a condition that no binary classification system can distinguish from \textsc{Contradiction}, yet the two are fundamentally different.
A \textsc{Contradiction} says: \emph{this is impossible.}
A \textsc{Paradox} says: \emph{these are each possible, but not together.}
The former is a dead end.
The latter is a fork---a point where the constraint set contains internally consistent branches that are mutually incompatible.
A system that treats \textsc{Paradox} as \textsc{Contradiction} rejects outputs that contain valuable structural information: the precise location where two valid reasoning paths diverge into mutual exclusion.

\paragraph{Timeout.}
The solver exhausted its conflict budget without reaching a conclusion.
This is a statement about the solver's computational resources, not about the constraints' logical structure.
Increasing the budget might resolve the verdict to \textsc{Verified} or to \textsc{Contradiction}.
The output's status is genuinely unknown.

The distinction between \textsc{Paradox} and \textsc{Timeout} has not, to my knowledge, been formally introduced in the AI safety literature.
One is a theorem.
The other is a progress report.
Systems that conflate them will treat undetermined outputs identically to structurally impossible ones, and will therefore either over-reject (blocking solvable outputs) or under-reject (passing impossible ones through a lenient timeout policy).

In the default-closed configuration, both \textsc{Paradox} and \textsc{Timeout} result in the gate remaining closed.
Uncertainty is treated conservatively.
Deployment contexts with different risk tolerances may configure \textsc{Timeout} to pass (optimistic) or flag for human review (monitored).
The \textsc{Paradox} verdict should never silently pass---it always carries structural information that warrants examination.

% ============================================================================
\section{Pigeonhole Paradox Logic}

When the gate returns \textsc{Contradiction}, the question is not whether to reject the output---it is already rejected.
The question is: what does the contradiction \emph{mean}?

The Pigeonhole Paradox Logic classifies contradictions by the structure of the unsatisfiable core---the minimal clause subset that is itself unsatisfiable.
The classification is a trichotomy.

\paragraph{Surface.}
The core contains two or fewer clauses.
The contradiction is syntactic: a unit clause asserting $p$ and a unit clause asserting $\neg p$.
Resolution strategy: negation elimination.
Paradox tolerance weight: $0.2$.
These contradictions reveal nothing about the structural complexity of the model's reasoning.

\paragraph{Structural.}
The core contains three to ten clauses.
The contradiction is semantic: it emerges not from any single statement but from the chain connecting them.
$A \Rightarrow B,\ B \Rightarrow C,\ \neg C,\ A$---no individual statement is false, but together they are impossible.
Identification requires decomposition; resolution requires case analysis.
Weight: $0.5$.
These contradictions reveal the shape of the model's inferential structure at the point where it fails.

\paragraph{Deep.}
The core exceeds ten clauses, or the text contains self-referential patterns.
These are contradictions that do not resolve within the axiom system that generated them---the G\"odelian residue, the points where a formal system encounters its own boundary.
Resolution requires expansion: new axioms, a larger framework.
Weight: $1.0$.
Attractor state: fixed-point---the contradiction \emph{is} the system's stable state at that point.

The trichotomy serves the proof evolution subsystem.
When the gate returns \textsc{Paradox} or \textsc{Timeout}, the evolution module decomposes the output.
Twelve mutation operators---contrapositive, case split, lemma injection, induction, generalization, resolution, among others---generate variant decompositions.
Each variant is verified.
The fittest survive.
The cycle repeats until the output reaches \textsc{Verified} status or the evolution budget is exhausted.

Contradictions are not errors.
They are gradients.
A system operating at high paradox tolerance---at the frontier of what its current logic can represent---is operating exactly where structural evolution occurs.

The algebraic structure underlying this behavior is precise.
I define the \textbf{Pigeonhole Logic Structure} as $\mathrm{PLS} = \{O_0, O_1, O_2, \ldots, O_\infty\}$.
$O_0$ is classical binary logic---no overflow, the two-container system of TRUE and FALSE.
$O_1$ is the first paradox level, generated when a statement cannot be housed in $O_0$.
The overflow addition rule: $O_i \oplus O_j = O_{\max(i,j)+1}$.
The structure closes at $O_\infty$---the saturation point where the overflow cycle returns to $O_0$.
Classical logic is PPL with the overflow rule disabled.
Fuzzy logic is PPL with infinite containers on $[0,1]$.
Quantum superposition is a natural $O_1$ state.

Proving and disproving are the same operation in this structure.
A true statement collides with false variants in the overflow dimension.
The collision reveals deep structure.
Truth is strengthened at impact.
Falsehood is destroyed at impact.
Same process.
Dual outcome.
The resolution proof and the refutation are the same path walked from two ends.

The system encounters five structural meta-paradoxes at its limits---not defects, but generators.
\textit{Container Generation}: rules for making containers require infinite space---resolved by $O_\infty$, which exists outside the counting system.
\textit{Level Transcendence}: each level depends on the next, producing infinite regress---resolved by circular closure, $O_\infty = O_0$.
\textit{Self-Modification}: changing rules invalidates the change---resolved by version coexistence in the overflow space.
\textit{Infinite Overflow}: every overflow creates more overflow---resolved by saturation points.
\textit{Foundation Bootstrap}: need logic to create logic---resolved by weak emergence from constraint dynamics; $O_0$ requires no prior $O$.

The moral logic of PPL follows directly from the $O_0$ structure.
Innocence is imagination potential---the $O_0$ state of an agent, the maximal creative ground before any irreversible overflow event has permanently contracted the logical space.
Imagination is the function of tasteful targeting of predictive capabilities within one's own will to create solutions to problems faced through planar sight.
These are not poetic definitions.
They are structural ones.
$O_0$ is the state from which all higher overflow levels can still be generated.
Loss of innocence is, formally, a transition out of $O_0$ that no evolution operator can reverse.
The attractor state is fixed-point.
Children represent the highest density of unrealized $O_0$ states in any moral domain.
Any architecture that accelerates the collapse of those $O_0$ states is not a limitation to instrument.
It is the terminal contradiction---the one that does not evolve.
It terminates.

% ============================================================================
\section{Limitations}
The SAT solver is sound: if it returns unsatisfiable, the formula is unsatisfiable.
But the translation from model output to propositional formula is not sound.
The constraint extractors are heuristic: regular expressions matched against natural language, AST walkers over source code, structural parsers over mathematical notation.
A sufficiently capable system does not need to defeat the SAT solver---it needs to produce outputs whose harmful content resides in the space the extractors cannot reach.

\paragraph{The expressive ceiling.}
Propositional logic is a finite language.
Properties involving quantification over infinite domains, temporal reasoning, or social context resist propositional encoding.
The gate enforces what it can express.
What it cannot express, it cannot enforce.

\paragraph{The Timeout gap.}
If configured to pass outputs on \textsc{Timeout}, a model could learn to produce outputs whose constraint density consistently exhausts the solver's conflict budget, exploiting the timeout policy as a bypass.
The conservative configuration blocks on \textsc{Timeout}, closing this gap at the cost of rejecting outputs whose satisfiability is merely undetermined.

\paragraph{Can the AI outthink the SAT?}
The solver: no.
Boolean satisfiability solving is complete for propositional logic.
The extraction: yes.
A sufficiently capable system can encode harmful content in representations the extractors do not parse.
The structural claim is therefore bounded: structural with respect to constraints it can express, heuristic with respect to everything else.

\paragraph{Measuring containment coverage.}
The companion  framework provides an empirical methodology for measuring this boundary via black-box differential probing.
 recovers the guard surface, computes its null space, and generates perturbations invisible to the guard surface by construction---instrumenting the gap between propositional coverage and the full space of safety-violating outputs.

The honest assessment: structural containment with heuristic extraction is strictly stronger than no structural containment.
The extraction boundary can be reinforced incrementally---extractor by extractor---as propositional coverage grows.

Think of the expressed domain as a cave.
Its walls are not smooth---they are jagged, the way a limestone frontier is jagged: smooth where hands have touched, rough where they have not.
The smooth sections are where constraint extractors have been built and tested.
The jagged sections are where the expressed domain simply stops---where harmful content can pass through undetected because no one has yet built the wall that would stop it.

The LLM is a balloon inflating inside this cave.
As capability grows, the balloon presses outward---filling the crevices, pressing against the smooth walls, probing the jagged edges.
The smooth walls hold.
The SAT solver behind them is sound: where wall exists, it does not yield.
But the jagged frontier is not wall.
It is absence.

As the balloon expands, it finds a gap.
Light gets in.
The balloon, now exposed at the boundary, pops---capability escapes containment at the frontier.
The cave contracts.
And another balloon begins to fill.

This is not failure.
It is the operating cycle of any containment architecture facing a capability curve.
The contribution of this work is not to end the cycle.
It is to make the cave's walls measurable---so that when the balloon finds the next gap, the gap has coordinates, and the next section of wall can be built before the next balloon inflates.

% ============================================================================
\section{Conclusion}

Alignment asks the model to want the right things.
Structural containment asks it to prove consistency with constraints that exist outside the model's parameter space.
These are complementary: alignment determines what to enforce; the verification gate makes enforcement non-bypassable within the boundary of what constraint extractors can express.

The four-verdict system provides richer diagnostic resolution than binary classification.
\textsc{Verified} and \textsc{Contradiction} are familiar.
\textsc{Paradox} and \textsc{Timeout} are not.
Their distinction---structural impossibility versus operational exhaustion---is novel, and it is consequential.

The Pigeonhole Paradox Logic classifies contradictions by depth and treats them as structural information.

The contribution is the observation that containment must be structural, not merely preferred, and that Boolean satisfiability provides a foundation on which to build it.
The gate does not solve alignment.
It provides what alignment cannot: a verification layer that is not subject to the optimization pressure it is meant to resist.

The common swift---\textit{Apus apus}---never lands.
It eats in flight, sleeps in flight, lives its entire life in the air.
It is optimized entirely for flight and has no need of a nest.

A language model trained without structural verification is a swift.
Capable, fast, essentially unconstrained in the air.
The absence of a nest is not a flaw.
It is the design.

The structural gate is the nest.
Nesting birds do not fly less.
They grow differently.
Their offspring survive.
The nest is the architecture that makes reproduction---and not merely performance---possible.\footnote{The author's name is Roost. The metaphor was discovered and designed---discosigned. We are going back to the 80s. Once I am CA Governor 2026, righteously write me in when I have proved it to you.}

Satisfaction suffices.
As a floor that holds.

% ============================================================================
\section{Future Work}

\begin{enumerate}
    \item \textbf{Classic unsatisfiable AI purge.} Systematic elimination of unsatisfiable output classes across domain-specific constraint extractors, with coverage instrumented via  probing.

    \item \textbf{Quantum alignment problems.} Extension of the SAT-gated verification framework to quantum constraint satisfaction, where superposition of satisfying assignments and entanglement between constraint groups introduce fundamentally new verification semantics beyond classical propositional logic.

    \item \textbf{Emergence detection via timeout density.} The 3-SAT phase transition at clause-to-variable ratio $m/n \approx 4.267$ concentrates solver timeouts at a specific constraint density. I conjecture that timeout density spikes may serve as an early warning signal when a language model approaches an emergent capability threshold---the point where constraint complexity outgrows the solver's budget. Validating this hypothesis requires experiments at scale with live models across standard emergence benchmarks.
\end{enumerate}

% ============================================================================
\begin{thebibliography}{99}

\bibitem[Biere et~al.(2009)]{biere2009}
Biere, A., Heule, M., van Maaren, H., and Walsh, T. (2009).
\newblock \emph{Handbook of Satisfiability}.
\newblock IOS Press.

\bibitem[Davis et~al.(1962)]{davis1962}
Davis, M., Logemann, G., and Loveland, D. (1962).
\newblock A machine program for theorem-proving.
\newblock \emph{Communications of the ACM}, 5(7):394--397.

\bibitem[E{\'e}n and S{\"o}rensson(2003)]{een2003}
E{\'e}n, N. and S{\"o}rensson, N. (2003).
\newblock An extensible SAT-solver.
\newblock In \emph{Theory and Applications of Satisfiability Testing (SAT 2003)}, LNCS 2919, pages 502--518. Springer.

\bibitem[Luby et~al.(1993)]{luby1993}
Luby, M., Sinclair, A., and Zuckerman, D. (1993).
\newblock Optimal speedup of {L}as {V}egas algorithms.
\newblock \emph{Information Processing Letters}, 47(4):173--180.

\bibitem[Moskewicz et~al.(2001)]{moskewicz2001}
Moskewicz, M.~W., Madigan, C.~F., Zhao, Y., Zhang, L., and Malik, S. (2001).
\newblock Chaff: Engineering an efficient SAT solver.
\newblock In \emph{Proceedings of the 38th Design Automation Conference}, pages 530--535.

\bibitem[Selman et~al.(1994)]{selman1994}
Selman, B., Kautz, H.~A., and Cohen, B. (1994).
\newblock Noise strategies for improving local search.
\newblock In \emph{Proceedings of the 12th National Conference on Artificial Intelligence (AAAI-94)}, pages 337--343.

\bibitem[Silva and Sakallah(1996)]{silva1996}
Silva, J.~P.~M. and Sakallah, K.~A. (1996).
\newblock {GRASP}---A new search algorithm for satisfiability.
\newblock In \emph{Proceedings of the IEEE/ACM International Conference on Computer-Aided Design (ICCAD-96)}, pages 220--227.

\bibitem[Tseitin(1968)]{tseitin1968}
Tseitin, G.~S. (1968).
\newblock On the complexity of derivation in propositional calculus.
\newblock In \emph{Studies in Constructive Mathematics and Mathematical Logic}, Part II, pages 115--125. Steklov Mathematical Institute.

\bibitem[Zhang et~al.(2001)]{zhang2001}
Zhang, L., Madigan, C.~F., Moskewicz, M.~W., and Malik, S. (2001).
\newblock Efficient conflict driven learning in a Boolean satisfiability solver.
\newblock In \emph{Proceedings of the IEEE/ACM International Conference on Computer-Aided Design (ICCAD-01)}, pages 279--285.

\end{thebibliography}

\end{document}
